%%%%%%%%%%%%%%%
% This CV example/template is based on my own
% CV which I (lamely attempted) to clean up, so that
% it's less of an eyesore and easier for others to use.
%
% LianTze Lim (liantze@gmail.com)
% 23 Oct, 2022
%
\documentclass[a4paper,skipsamekey,11pt,english]{curve}

% Uncomment to enable Chinese; needs XeLaTeX
% \usepackage{ctex}


% Default biblatex style used for the publication list is APA6. If you wish to use a different style or pass other options to biblatex you can change them here. 
\PassOptionsToPackage{style=ieee,sorting=none,uniquename=init,defernumbers=true,maxbibnames=30}{biblatex}

% Most commands and style definitions are in settings.sty.
\usepackage{settings}

% If you need to further customise your biblatex setup e.g. with \DeclareFieldFormat etc please add them here AFTER loading settings.sty. For example, to remove the default "[Online] Available:" prefix before URLs when using the IEEE style:
\DefineBibliographyStrings{english}{url={\textsc{url}}}

%% Publications: all .bib files go into one bibliography, so that entries can be
%% cited anywhere in the CV (e.g. \cite{Mungiello2026} -> [J3], clickable).
%% Each file gets a keyword (used to split the lists) and a label prefix.
\addbibresource{journal_papers.bib}
\addbibresource{journal_papers_under_review.bib}
\addbibresource{conference_papers.bib}
\addbibresource{conference_papers_under_review.bib}
\addbibresource{workshop_papers.bib}
\addbibresource{other_publications.bib}
\addbibresource{poster_sessions.bib}
\addbibresource{patents.bib}
\DeclareSourcemap{\maps[datatype=bibtex, overwrite=true]{
  \map{\perdatasource{journal_papers.bib}\step[fieldset=keywords, fieldvalue={,}, appendstrict]\step[fieldset=keywords, fieldvalue={cvjournal}, append]}
  \map{\perdatasource{journal_papers_under_review.bib}\step[fieldset=keywords, fieldvalue={,}, appendstrict]\step[fieldset=keywords, fieldvalue={cvjournalrev}, append]}
  \map{\perdatasource{conference_papers.bib}\step[fieldset=keywords, fieldvalue={,}, appendstrict]\step[fieldset=keywords, fieldvalue={cvconf}, append]}
  \map{\perdatasource{conference_papers_under_review.bib}\step[fieldset=keywords, fieldvalue={,}, appendstrict]\step[fieldset=keywords, fieldvalue={cvconfrev}, append]}
  \map{\perdatasource{workshop_papers.bib}\step[fieldset=keywords, fieldvalue={,}, appendstrict]\step[fieldset=keywords, fieldvalue={cvworkshop}, append]}
  \map{\perdatasource{other_publications.bib}\step[fieldset=keywords, fieldvalue={,}, appendstrict]\step[fieldset=keywords, fieldvalue={cvother}, append]}
  \map{\perdatasource{poster_sessions.bib}\step[fieldset=keywords, fieldvalue={,}, appendstrict]\step[fieldset=keywords, fieldvalue={cvposter}, append]}
  \map{\perdatasource{patents.bib}\step[fieldset=keywords, fieldvalue={,}, appendstrict]\step[fieldset=keywords, fieldvalue={cvpatent}, append]}
}}
\assignrefcontextkeyws[labelprefix=J]{cvjournal}
\assignrefcontextkeyws[labelprefix=JR]{cvjournalrev}
\assignrefcontextkeyws[labelprefix=C]{cvconf}
\assignrefcontextkeyws[labelprefix=CR]{cvconfrev}
\assignrefcontextkeyws[labelprefix=W]{cvworkshop}
\assignrefcontextkeyws[labelprefix=O]{cvother}
\assignrefcontextkeyws[labelprefix=PS]{cvposter}
\assignrefcontextkeyws[labelprefix=PT]{cvpatent}

%% Reverse numbering within each list (first journal paper = J<total>), for both
%% list labels and citations. Totals are counted while printing and stored in
%% the .aux file, so they are applied from the second LaTeX run on.
\makeatletter
\AtEveryBibitem{%
  \edef\cv@prefix{\thefield{labelprefix}}%
  \ifcsdef{cv@n@\cv@prefix}{}{\csgdef{cv@n@\cv@prefix}{0}}%
  \csxdef{cv@n@\cv@prefix}{\the\numexpr\csuse{cv@n@\cv@prefix}+1\relax}}
\newcommand*{\cv@writetotal}[1]{%
  \ifcsdef{cv@n@#1}
    {\immediate\write\@auxout{\gdef\string\cv@total@#1{\csuse{cv@n@#1}}}}{}}
\AtEndDocument{\forcsvlist{\cv@writetotal}{J,JR,C,CR,W,O,PS,PT}}
\DeclareFieldFormat{labelnumber}{%
  \ifcsdef{cv@total@\thefield{labelprefix}}
    {\the\numexpr\csuse{cv@total@\thefield{labelprefix}}+1-#1\relax}
    {#1}}
\makeatother

%% Specify your last name(s) and first name(s) (as given in the .bib) to automatically bold your own name in the publications list. 
%% One caveat: You need to write \bibnamedelima where there's a space in your name for this to work properly; or write \bibnamedelimi if you use initials in the .bib
% \mynames{Lim/Lian\bibnamedelima Tze}

%% You can specify multiple names like this, especially if you have changed your name or if you need to highlight multiple authors. See items 6–9 in the example "Journal Articles" output.
\mynames{Piccinini/Mattia}
%% MAKE SURE THERE IS NO SPACE AFTER THE FINAL NAME IN YOUR \mynames LIST

% Change the fonts if you want
\ifxetexorluatex % If you're using XeLaTeX or LuaLaTeX
  \usepackage{fontspec} 
  %% You can use \setmainfont etc; I'm just using these font packages here because they provide OpenType fonts for use by XeLaTeX/LuaLaTeX anyway
  \usepackage[p,osf,swashQ]{cochineal}
  \usepackage[medium,bold]{cabin}
  \usepackage[varqu,varl,scale=0.9]{zi4}
\else % If you're using pdfLaTeX or latex
  \usepackage[T1]{fontenc}
  \usepackage[p,osf,swashQ]{cochineal}
  \usepackage{cabin}
  \usepackage[varqu,varl,scale=0.9]{zi4}
\fi


% Change the page margins if you want
% \geometry{left=1cm,right=1cm,top=1.5cm,bottom=1.5cm}

% Change the colours if you want
% \definecolor{SwishLineColour}{HTML}{00FFFF}
% \definecolor{MarkerColour}{HTML}{0000CC}

% Change the item prefix marker if you want
% \prefixmarker{$\diamond$}

%% Photo is only shown if "fullonly" is included
\includecomment{fullonly}
% \excludecomment{fullonly}

%%%%%%%%%%%%%%%%%%%%%%%%%%%%%%%%%%%%%%

\leftheader{%
  \begin{center}
        \LARGE\bfseries\sffamily Mattia Piccinini
  \end{center}
  \vspace{-0.8cm}
  \begin{center}
        \large Humboldt Post-doctoral Fellow\\Technical University of Munich (TUM), Germany
  \end{center}

  \begin{tabular}{p{0.22\textwidth} p{0.28\textwidth} p{0.22\textwidth}}
     %\makefield{\faMapPin}{Munich, Germany} &
     % create a website field
      \makefield{\faGlobe}{\href{https://mattiapiccinini95.github.io}{Website}} & 
     \makefield{\faEnvelope[regular]}{\href{mailto:mattia.picci95@gmail.com}{mattia.picci95@gmail.com}} & \makefield{\faPhone}{+39 3382383705} \\ 
     \makefield{\faLinkedin}
     {\href{https://www.linkedin.com/in/mattia-piccinini}{LinkedIn}} & 
     \makefield{\faResearchgate}{\href{https://www.researchgate.net/profile/Mattia-Piccinini}{Research Gate}}
     & \makefield{\faGraduationCap}{\href{https://scholar.google.com/citations?user=xtgKwEoAAAAJ\&hl=it\&authuser=1}{Google Scholar}} 
     %\makefield{\aiGoogleScholar {\href{https://scholar.google.com/citations?user=xtgKwEoAAAAJ\&hl=it\&authuser=1}{Mattia Piccinini}}
  \end{tabular}
}

\rightheader{~}
\begin{fullonly}
\photo[r]{figures/Foto_2024.JPG}
\photoscale{0.15}
\end{fullonly}

\title{Curriculum Vitae}

\begin{document}
\nocite{*}
\makeheaders[c]

\makerubrichead{Research Interests}
%
My research develops autonomous \textbf{robotic systems} that can reason, plan, and act in uncertain real-world environments, under physical constraints and safety requirements. While AI has enabled robots to learn complex tasks from data, purely data-driven approaches often lack interpretability, require large datasets, and fail to guarantee safe, reliable behavior. I aim to bridge this gap by combining \textbf{physics-informed} knowledge and structured \textbf{optimization} principles with \textbf{machine learning}, producing planning and control methods that are \textbf{data-efficient}, adaptive, and safety-aware. I couple methodological advances with \textbf{real-world} validation on robotic platforms, targeting practical impact on the next-generation autonomous systems.
%
% The goal of my research is to develop autonomous robotic systems that can reason, plan, and act in uncertain real-world environments, while respecting complex physical constraints, safety requirements, and computational limitations. Recent advances in artificial intelligence have shown how robots can learn from data to perform complex tasks. However, purely data-driven approaches require massive datasets to generalize, lack interpretability, and often fail to provide safe and reliable performance, which is critical for real-world robotics. My research aims to bridge these gaps by exploiting prior physical knowledge and structured optimization principles to develop physics-informed, learning-based robot planning and control methods that are data-efficient, adaptive in changing environments, and compatible with real-time constraints. I combine methodological advances with real-world validation on robotic platforms, aiming for practical impact on the future autonomous systems.
%
% My research lies at the intersection of \textbf{robotics}, \textbf{machine learning}, and \textbf{control} engineering. I develop algorithms for real-time trajectory planning, control, and state estimation, with a focus on mobile ground robots in uncertain, dynamic environments. I like to integrate \textbf{prior knowledge} of the robot dynamics into learning-based methods, such as neural networks, to enhance \textbf{generalization} to unseen scenarios with minimal training data.

\vspace{0.4cm}

% \begin{itemize}
%     \item Trajectory planning and control with autonomous vehicles \vspace{-0.2cm}
%     \item Real-time optimization and model predictive control \vspace{-0.2cm}
%     \item Model-structured neural networks for modeling, control and estimation of dynamical systems \vspace{-0.2cm}
%     \item Learning to control and system identification
% \end{itemize}

\makerubric{blocks/education}
\makerubric{blocks/academic_experience}
\makerubric{blocks/teaching_experience}
\makerubric{blocks/work_experience}
%% Sometimes when a section can't be nicely modelled with the \entry[]... mechanism; hack our own
\makerubrichead{Publications and Patents}

\vspace{0.1cm}
\begin{center}
\begin{tabular}{|l|l|}
\hline
\textbf{Metric} & \textbf{Value} \\
\hline
N. journal papers & $15$ \\
N. conference papers & $23$ \\
N. patents & $2$ ($1$ pending) \\
N. citations (Google Scholar) & $530+$ \\
h-index (Google Scholar) & $13$ \\
\hline
\end{tabular}
\end{center}

% #1: label prefix, #2: keyword assigned to the source .bib file (see main_cv.tex), #3: title
\newcommand{\printpublist}[3]{%
	\newrefcontext[labelprefix=#1]%
	\printbibliography[keyword=#2,heading={subbibliography},title={#3},resetnumbers=true]%
}

\vspace{0.1cm}

\printpublist{J}{cvjournal}{Peer-Reviewed Journal Papers}
\printpublist{JR}{cvjournalrev}{Peer-Reviewed Journal Papers \textit{Under Review}}

\printpublist{C}{cvconf}{Peer-Reviewed Conference Papers}
% \printpublist{CR}{cvconfrev}{Peer-Reviewed Conference Papers \textit{Under Review}}

\printpublist{W}{cvworkshop}{Peer-Reviewed Workshop Papers}
\printpublist{O}{cvother}{Other Publications}
% \printbibliographyfromfile{invited_talks.bib}{Invited Talks}
\printpublist{PS}{cvposter}{Poster Sessions}

\printpublist{PT}{cvpatent}{Patents}
\endrefcontext

\makerubric{blocks/awards}
\makerubric{blocks/projects}
\makerubrichead{Academic Engagement}

{\large\bfseries\sffamily Editorial Activities}

\begin{itemize}
    \item \textbf{Associate Editor}, IEEE/RSJ International Conference on Intelligent Robots and Systems (IROS), 2026
    \item \textbf{Associate Editor}, IEEE International Conference on Intelligent Transportation Systems (ITSC), 2026
    \item \textbf{Invited Session Organizer}, Invited Session on ``Automated Vehicles at the Limits: Lessons and Opportunities from the Racetrack'', IEEE International Conference on Intelligent Transportation Systems (ITSC), 2026
    \item \textbf{Guest Editor} of the Special Issue on ``Methods and Algorithms for Autonomous Systems in Intelligent Transportation'', IEEE Open Journal of Intelligent Transportation Systems, from September 2026
\end{itemize}

\vspace{0.2cm}
{\large\bfseries\sffamily Conference Committees}

\begin{itemize}
    \item \textbf{Chair} of the session ``Automated Vehicles at the Limits: Lessons and Opportunities from the Racetrack'', at the IEEE International Conference on Intelligent Transportation Systems (ITSC), Naples (Italy), 2026
    \item \textbf{Chair} of the session ``Tracking and Coordinating Autonomous Vehicle Fleets'', at the IEEE/RSJ International Conference on Intelligent Robots and Systems (IROS), Pittsburgh (USA), 2026
    \item \textbf{Chair} of the session ``Model Predictive Control'', at the Automatica.it conference, Parma (Italy), 2026
    \item \textbf{Chair} of the session ``Path Planning'', at the Automatica.it conference, Bolzano (Italy), 2024
    \item \textbf{Co-chair} of the session ``Modeling, Control, and Optimization of Transportation Systems'', at the 17th IFAC Symposium on Control in Transportation Systems (CTS), Ayia Napa (Cyprus), 2024
\end{itemize}

\vspace{0.2cm}
{\large\bfseries\sffamily Workshop Organization}

\begin{itemize}
    \item \textbf{Workshop Organizer}: M. Piccinini, M. Salazar, J. Betz, ``Automated and Electrified Racing: Opportunities for Future Intelligent Vehicles'', IEEE International Conference on Intelligent Transportation Systems (ITSC), 2026
    \item \textbf{Workshop Organizer}: D. Meli, M. Piccinini, F. Trotti, A. Farinelli, J. Betz, ``AI Meets Control for Resilient and Reliable Robot Planning (AIM-Ctrl)'', IEEE/RSJ International Conference on Intelligent Robots and Systems (IROS), 2026
\end{itemize}

% \vspace{0.2cm}
% The following workshop proposals are currently \textit{under review}:
% \begin{itemize}
%     \item \textbf{Workshop Organizer}: M. Piccinini, A. Plebe, B. Zarrouki, D. Wang, G. P. Rosati Papini, J. Betz, ``Data Meets Physics: Learning with Physical Priors in Robotics'', IEEE/RSJ International Conference on Intelligent Robots and Systems (IROS), 2026
% \end{itemize}

% \vspace{0.2cm}
{\large\bfseries\sffamily Invited Talks}

\begin{enumerate}
    \item Invited talk at the 16th Workshop on Planning, Perception and Navigation for Intelligent Vehicles, 2026 IEEE/RSJ International Conference on Intelligent Robots and Systems (IROS), Pittsburgh, USA: \textit{Differentiable Physics-Informed Learning and Control in Autonomous Vehicle Racing}, 2026
    \item Invited talk at the Workshop ``Automated Mobility in Mixed Traffic - Toward Robust, Human-Centric, and Trustworthy Systems'', IEEE 29th International Conference on Intelligent Transportation Systems (ITSC), Naples, Italy: \textit{From the Lab to the Street: Open-Source Software for Interactive Urban Scenario Driving}, 2026
    \item Invited talk at the Workshop ``An Introduction to Autoware and Its Application Platforms'', IEEE 29th International Conference on Intelligent Transportation Systems (ITSC), Naples, Italy: \textit{Open-Source Autonomy: A Path from Interactive Simulation to Real-World Driving}, 2026
    \item Invited talk at the University of Verona: \textit{Physics-Encoded Trajectory Planning and Control: From Racing Cars to General Robotics}, 2026
    \item Invited talk at Eindhoven University of Technology: \textit{Physics-Guided Motion Generation and Control: From Autonomous Racing Towards General Robotics}, 2025
    \item Invited talk at the online Workshop series organized by the Ekumen company: \textit{NNodely: An Open Framework for Model-Structured Neural Networks in Robotics}, 2025, \url{https://www.youtube.com/watch?v=mz70WepnwMM}
    \item Invited talk at the Technical University of Munich: \textit{Artificial Drivers to Learn the Vehicle Dynamics, Plan and Execute Online Time-Optimal Maneuvers}, 2024
    \item Invited talk at the University of Trento, Seminar for PhD students: \textit{Model-Structured Neural Networks to Model and Control Physical Systems}, 2024
    \item Invited talk at Universität der Bundeswehr, Munich: \textit{Virtual race drivers for planning and control with robotic cars}, 2022
\end{enumerate}

\vspace{0.2cm}
{\large\bfseries\sffamily University Service}

\begin{itemize}
    \item \textbf{TUM Internationalization} support: organized the research visits of Prof. Basilio Lenzo, Prof. Francesco Biral, Dr. Mattia Laurini, and the visiting Ph.D. students Aniello Mungiello, Gonçalo Torres, Mattia Piazza, and Gianluca Frison to the Technical University of Munich, Professorship of Autonomous Vehicle Systems (2025 -- present).
    \item \textbf{Final Examination Committee} Member for around $40$ Master's theses at the Technical University of Munich, Professorship of Autonomous Vehicle Systems (2024 -- present).
    \item \textbf{Final Examination Committee} Member for around $20$ Bachelor's theses at the Technical University of Munich, Professorship of Autonomous Vehicle Systems (2024 -- present).
    \item \textbf{Selection Committee Chair} for external visiting Ph.D. positions and M.Sc. theses at the Technical University of Munich, Professorship of Autonomous Vehicle Systems (2025 -- present).
    \item \textbf{Open Days} Lab Presenter, Master of Science in Mechatronics Engineering, University of Trento (2020 -- 2022).
\end{itemize}

\vspace{0.2cm}
{\large\bfseries\sffamily Society Memberships}

\begin{itemize}
    \item IEEE Robotics and Automation Society (RAS)
    \item IEEE Intelligent Transportation Systems Society (ITSS)
    \item Società Italiana dei Docenti e Ricercatori di Automatica (SIDRA)
    \item German Robotics Society (Deutsche Gesellschaft für Robotik, DGR)
\end{itemize}

\vspace{0.2cm}
{\large\bfseries\sffamily Attendance at Conferences with Oral Presentations}

\begin{itemize}
    \item IEEE International Conference on Intelligent Transportation Systems (ITSC) 2026, Naples (Italy)
    \item German Robotics Conference 2026, Cologne (Germany)
    \item Automatica.it 2026, Parma (Italy)
    \item IEEE International Conference on Robotics and Automation (ICRA) 2025, Atlanta (USA)
    \item IEEE/RSJ International Conference on Intelligent Robots and Systems (IROS) 2025, Hangzhou (China)
    \item Automatica.it 2025, Perugia (Italy)
    \item IFAC Symposium on Control in Transportation Systems (CTS) 2024, Ayia Napa (Cyprus)
    \item Automatica.it 2024, Bolzano (Italy)
    \item IEEE/RSJ International Conference on Intelligent Robots and Systems (IROS) 2020, Las Vegas (USA) 
\end{itemize}

\vspace{0.2cm}
{\large\bfseries\sffamily Peer-Review Activities}

\vspace{0.3cm}
{\normalsize\bfseries\sffamily Journals}: 
IEEE Transactions on Intelligent Transportation Systems; IEEE Transactions on Intelligent Vehicles; IEEE Transactions on Vehicular Technology; IEEE Transactions on Mechatronics; IEEE Transactions on Industrial Electronics; IEEE Open Journal of Control Systems; IEEE Transactions on Control Systems Technology; IEEE Transactions on Robotics; IEEE Robotics and Automation Letters (RA-L); IEEE Access; Robotics and Autonomous Systems; Springer Nature - International Journal of Machine Learning and Cybernetics; Springer Optimization and Engineering; Springer Nature Communications; Control Engineering Practice (IFAC); Vehicle System Dynamics; Mechanical Systems and Signal Processing; Results in Materials (Elsevier); MDPI Vehicles; MDPI Machines; MDPI Electronics; MDPI Actuators; MDPI Algorithms; MDPI Drones; MDPI Symmetry; MDPI Applied Sciences; MDPI World Electric Vehicle Journal 

\vspace{0.3cm}
{\normalsize\bfseries\sffamily Conferences}: IEEE International Conference on Robotics and Automation (ICRA); IEEE/RSJ International Conference on Intelligent Robots and Systems (IROS); IEEE Intelligent Vehicles Symposium (IV); IEEE International Conference on Intelligent Transportation Systems (ITSC); IEEE Conference on Decision and Control (CDC); Conference on Neural Information Processing Systems (NeurIPS); European Control Conference (ECC); IEEE Conference on Control Technology and Applications; IEEE International Conference on Real-time Computing and Robotics; IEEE International Conference on Advanced Motion Control; International Conference on System Theory, Control and Computing (ICSTCC); IFAC Symposium on Control in Transportation Systems; Winter Conference on Applications of Computer Vision (WACV).
\vspace{0.3cm}
\makerubric{blocks/skills}
\newpage
\makerubrichead{Personal Details and Interests}
\begin{itemize}
    \item Date of birth: 16/07/1995 \vspace{-0.2cm}
    \item Nationality: Italian \vspace{-0.2cm}
    \item Personal interests: traveling, swimming, hiking 
\end{itemize}
\makerubrichead{References}
\begin{itemize}
    \item Prof. Johannes Betz, Technische Universit\"at M\"unchen, M\"unchen (Germany), \\ phone: +49 89 289 10407, email: \href{mailto:johannes.betz@tum.de}{johannes.betz@tum.de} \vspace{-0.2cm}
	\item Prof. Francesco Biral, University of Trento, Trento (Italy), \\ phone: +39 (0)461 282513, email: \href{mailto:francesco.biral@unitn.it}{francesco.biral@unitn.it} \vspace{-0.2cm}
	\item Prof. Matthias Gerdts, Universit\"at der Bundeswehr, M\"unchen (Germany), \\ phone: +49 (0)89 6004 2672, email: \href{mailto:matthias.gerdts@unibw.de}{matthias.gerdts@unibw.de} \vspace{-0.2cm}
	\item Prof. Mauro Da Lio, University of Trento, Trento (Italy), \\ phone: +39 (0)461 282501, email: \href{mailto:mauro.dalio@unitn.it}{mauro.dalio@unitn.it} \vspace{-0.2cm}
    \item Prof. Gastone Pietro Rosati Papini, University of Trento, Trento (Italy), \\ phone: +39 (0)461 285352, email: \href{mailto:gastone.rosatipapini@unitn.it}{gastone.rosatipapini@unitn.it}
\end{itemize}
\vspace{1cm}
\noindent \textit{Munich, 25/09/2026}

\vspace{0.2cm}
\begin{flushright}
\begin{tabular}{m{5cm}}
\\ \hline
\centering Mattia Piccinini \\
\end{tabular}
\end{flushright}

\end{document}